\section{Research Objectives and Design Principles}
The audit operationalizes six research questions. \textbf{RQ1} asks which source and analytical components are actually integrated and where manual interfaces remain. \textbf{RQ2} asks which provenance invariants are enforced and what they fail to establish. \textbf{RQ3} asks whether similarity and conditional-randomization outputs reproduce on frozen inputs. \textbf{RQ4} concerns separation of communication timestamps, supplied physical-event intervals, and unknown time. \textbf{RQ5} asks which semantic failure modes survive software-level quality gates. \textbf{RQ6} asks which bounded analyst workflow is supported without implying independently validated investigation.

Three design principles organize the implementation. \emph{Typed separation} keeps sources, artifacts, statements, evidence assessments, and findings distinct. \emph{Explicit incompleteness} retains missing inputs, failed collection, unknown lineage, unresolved events, and unverified claims. \emph{Inspectable computation} records inputs, parameters, versions, and outputs so that an analyst can examine how a result arose. These principles are partly implemented and partly limited by legacy interfaces; they are not asserted as universal invariants across every script.

Our method is an implementation audit with bounded empirical checks. We prioritize current code and fresh computations over completion reports and plans. A successful test establishes only its tested property. An output count establishes inventory, not correctness. A curated example demonstrates a workflow, not blind accuracy. The new publication artifacts preserve these distinctions in a claim--evidence matrix and machine-readable experiment manifest. No core Social implementation or previous case-study manuscript was changed to improve results reported here.
